\section{Review of Measure Theory}

\subsection{Construction of measures}

\begin{de}[Rings, semirings, and semialgebras]
Let $\Omega$ be a set.
\begin{enumerate}[label=(\roman*)]
  \item A family $\mathcal R\subseteq\mathcal P(\Omega)$ is a \emph{ring} if
  $\varnothing\in\mathcal R$ and $A\cup B,A\setminus B\in\mathcal R$
  whenever $A,B\in\mathcal R$. It is an \emph{algebra} if, in addition,
  $\Omega\in\mathcal R$.
  \item A family $\mathcal E\subseteq\mathcal P(\Omega)$ is a
  \emph{semiring} if $\varnothing\in\mathcal E$, it is closed under finite
  intersections, and every difference $A\setminus B$ with $A,B\in\mathcal E$
  is a finite disjoint union of members of $\mathcal E$.
  \item A semiring containing $\Omega$ is a \emph{semialgebra}.
\end{enumerate}
\end{de}

\begin{thm}[Ring generated by a semiring]
If $\mathcal E$ is a semiring on $\Omega$, then the ring generated by
$\mathcal E$ is
\[
\mathcal R(\mathcal E)
 =\left\{\bigsqcup_{j=1}^{m}E_j:
 m\in\mathbb N_0,\ E_j\in\mathcal E\right\},
\]
where the empty disjoint union is $\varnothing$. If $\mathcal E$ is a
semialgebra, this ring contains $\Omega$ and is therefore the algebra generated
by $\mathcal E$.
\end{thm}

\begin{proof}
Let $\mathcal R_0$ denote the displayed family. It contains $\mathcal E$.
Intersections of two members are finite disjoint unions of sets
$E_i\cap F_j\in\mathcal E$. Repeatedly using the semiring difference property
shows that
\[
E_i\setminus\bigcup_jF_j
\]
is a finite disjoint union of members of $\mathcal E$; taking the union over
$i$ proves closure under differences. Hence $\mathcal R_0$ is a ring. Every
ring containing $\mathcal E$ contains all finite disjoint unions of its
members, so $\mathcal R_0$ is the smallest such ring. The final assertion is
immediate when $\Omega\in\mathcal E$.
\end{proof}

\begin{de}[Premeasure]
Let $\mathcal R$ be a ring on $\Omega$. A map
$\mu_0:\mathcal R\to[0,\infty]$ is a \emph{premeasure} if
$\mu_0(\varnothing)=0$ and, whenever $(A_n)$ is a disjoint sequence in
$\mathcal R$ whose union also belongs to $\mathcal R$,
\[
\mu_0\!\left(\bigsqcup_{n=1}^{\infty}A_n\right)
 =\sum_{n=1}^{\infty}\mu_0(A_n).
\]
The same definition is used on a semiring, with every set in the displayed
identity required to belong to the semiring.
\end{de}

\begin{thm}[Extension from a semiring to its generated ring]
Let $\mathcal E$ be a semiring and let
$\chi:\mathcal E\to[0,\infty]$ be finitely additive in the following sense:
if $E=\bigsqcup_{j=1}^{m}E_j$ with all sets in $\mathcal E$, then
$\chi(E)=\sum_{j=1}^{m}\chi(E_j)$. For
$A=\bigsqcup_{j=1}^{m}E_j\in\mathcal R(\mathcal E)$, define
\[
\widetilde\chi(A):=\sum_{j=1}^{m}\chi(E_j).
\]
Then $\widetilde\chi$ is well defined and finitely additive. If $\chi$ is a
premeasure on $\mathcal E$, then $\widetilde\chi$ is a premeasure on
$\mathcal R(\mathcal E)$.
\end{thm}

\begin{proof}
Suppose
$A=\bigsqcup_{i=1}^{m}E_i=\bigsqcup_{j=1}^{r}F_j$. The sets
$E_i\cap F_j$ form a common finite disjoint refinement. Finite additivity on
$\mathcal E$ gives
\[
\sum_i\chi(E_i)
 =\sum_{i,j}\chi(E_i\cap F_j)
 =\sum_j\chi(F_j),
\]
so the definition is independent of the representation. Finite additivity of
$\widetilde\chi$ follows by joining disjoint representations.

Now assume that $\chi$ is a premeasure. Let
$A=\bigsqcup_{n\ge1}A_n$ with $A,A_n\in\mathcal R(\mathcal E)$. Refine $A$
and every $A_n$ into semiring pieces. Each semiring piece of $A$ is a
countable disjoint union of its intersections with the pieces of the $A_n$.
Countable additivity of $\chi$ on $\mathcal E$, followed by Tonelli's theorem
for nonnegative series, yields
$\widetilde\chi(A)=\sum_n\widetilde\chi(A_n)$.
\end{proof}

\begin{prop}[Useful criterion for premeasures]
Let $\mu_0$ be finitely additive on a ring $\mathcal R$. Then $\mu_0$ is a
premeasure if and only if it is countably subadditive on $\mathcal R$:
whenever $A\subseteq\bigcup_{n\ge1}A_n$ with $A,A_n\in\mathcal R$,
\[
\mu_0(A)\leq\sum_{n=1}^{\infty}\mu_0(A_n).
\]
The same criterion holds on a semiring after passing to its generated ring by
the preceding theorem.
\end{prop}

\begin{proof}
A premeasure is countably subadditive after disjointifying the covering sets
and intersecting them with $A$. Conversely, if
$A=\bigsqcup_nA_n\in\mathcal R$, countable subadditivity gives
$\mu_0(A)\leq\sum_n\mu_0(A_n)$. Finite additivity and monotonicity give, for
each $N$,
$\mu_0(A)\geq\sum_{n=1}^{N}\mu_0(A_n)$. Letting $N\to\infty$ gives the
reverse inequality.
\end{proof}

\begin{thm}[Lebesgue--Stieltjes premeasure]
Let $F:\mathbb R\to\mathbb R$ be nondecreasing. On the semiring
\[
\mathcal I=\{\varnothing\}\cup\{(a,b]:a,b\in\mathbb R,
 a<b\},
\]
define
\[
\chi_F((a,b])=F(b)-F(a),\qquad \chi_F(\varnothing)=0.
\]
Then $\chi_F$ is a premeasure if and only if $F$ is right-continuous. In that
case its Carath\'eodory extension is the Lebesgue--Stieltjes measure associated
with $F$.
\end{thm}

\begin{proof}
Finite additivity follows by telescoping after sorting endpoints. Assume first
that $F$ is right-continuous. By the preceding criterion, it is enough to
show countable subadditivity. Suppose
$(a,b]\subseteq\bigcup_n(a_n,b_n]$. Fix $\varepsilon>0$. Choose
$0<\delta<b-a$ so that
$F(a+\delta)-F(a)<\varepsilon$, and choose $\eta_n>0$ so that
$F(b_n+\eta_n)-F(b_n)<\varepsilon2^{-n}$. The open intervals
$(a_n,b_n+\eta_n)$ cover the compact interval $[a+\delta,b]$, so a finite
subfamily covers it. The elementary finite interval-cover inequality for a
nondecreasing function gives
\[
F(b)-F(a+\delta)
 \leq\sum_{n\in J}\bigl(F(b_n+\eta_n)-F(a_n)\bigr)
\]
for that finite index set $J$. Therefore
\[
F(b)-F(a)
 \leq\sum_{n=1}^{\infty}\bigl(F(b_n)-F(a_n)\bigr)+2\varepsilon.
\]
Let $\varepsilon\downarrow0$.

Conversely, suppose $\chi_F$ is a premeasure and let $h_n\downarrow0$.
The intervals $(x+h_{n+1},x+h_n]$ are disjoint and their union is
$(x,x+h_1]$. Countable additivity gives
\[
F(x+h_1)-F(x)
 =\sum_{n=1}^{\infty}
   \bigl(F(x+h_n)-F(x+h_{n+1})\bigr).
\]
The tail of this convergent nonnegative series is
$F(x+h_n)-F(x)$, which therefore tends to zero. Thus $F$ is
right-continuous.
\end{proof}

\begin{de}[Outer measure]
An \emph{outer measure} on $\Omega$ is a map
$\mu^*:\mathcal P(\Omega)\to[0,\infty]$ such that
\begin{enumerate}[label=(\roman*)]
  \item $\mu^*(\varnothing)=0$;
  \item $A\subseteq B$ implies $\mu^*(A)\leq\mu^*(B)$;
  \item $\mu^*(\bigcup_nA_n)\leq\sum_n\mu^*(A_n)$.
\end{enumerate}
\end{de}

\begin{thm}[Outer measure generated by a covering cost]
Let $\mathcal E\subseteq\mathcal P(\Omega)$ contain $\varnothing$, and let
$\chi:\mathcal E\to[0,\infty]$ satisfy $\chi(\varnothing)=0$. With the
convention $\inf\varnothing=\infty$, define
\[
\chi^*(A)=
\inf\left\{\sum_{n=1}^{\infty}\chi(E_n):
 A\subseteq\bigcup_{n=1}^{\infty}E_n,
 E_n\in\mathcal E\right\}.
\]
Then $\chi^*$ is an outer measure.
\end{thm}

\begin{proof}
The empty set is covered by the empty set, so $\chi^*(\varnothing)=0$.
Monotonicity is immediate because every cover of $B$ also covers every
$A\subseteq B$. For countable subadditivity, fix $\varepsilon>0$ and, for
each $k$, choose a cover $(E_{k,j})_j$ of $A_k$ with total cost at most
$\chi^*(A_k)+\varepsilon2^{-k}$ whenever $\chi^*(A_k)<\infty$. The combined
family covers $\bigcup_kA_k$. Taking the infimum and then letting
$\varepsilon\downarrow0$ proves the claim; cases with infinite right-hand
side are automatic.
\end{proof}

\begin{thm}[Agreement with a premeasure]
Let $\mu_0$ be a premeasure on a ring $\mathcal R$, and let $\mu^*$ be
the outer measure obtained by covering with sets in $\mathcal R$. Then
\[
\mu^*(A)=\mu_0(A)\qquad(A\in\mathcal R).
\]
\end{thm}

\begin{proof}
The one-set cover gives $\mu^*(A)\leq\mu_0(A)$. Conversely, if
$A\subseteq\bigcup_nA_n$ with $A_n\in\mathcal R$, disjointify the cover and
intersect it with $A$. Countable additivity of $\mu_0$ then gives
$\mu_0(A)\leq\sum_n\mu_0(A_n)$. Taking the infimum over all covers proves the
reverse inequality.
\end{proof}

\begin{de}[Carath\'eodory measurability]
For an outer measure $\mu^*$, a set $E\subseteq\Omega$ is
\emph{$\mu^*$-measurable} if
\[
\mu^*(A)=\mu^*(A\cap E)+\mu^*(A\setminus E)
\qquad\text{for every }A\subseteq\Omega.
\]
\end{de}

\begin{thm}[Carath\'eodory's theorem]
The family $\mathcal M_{\mu^*}$ of $\mu^*$-measurable sets is a
$\sigma$-algebra, and $\mu=\mu^*|_{\mathcal M_{\mu^*}}$ is a complete
measure. If $\mu^*$ is generated by a premeasure $\mu_0$ on a ring
$\mathcal R$, then
\[
\mathcal R\subseteq\mathcal M_{\mu^*}
\quad\text{and}\quad
\mu|_{\mathcal R}=\mu_0.
\]
\end{thm}

\begin{proof}
The reverse inequality in the defining equality follows from subadditivity,
so only the forward inequality must be checked. The defining condition is
stable under complements. Applying it successively shows that it is stable
under finite unions. If $(E_n)$ is disjoint and measurable, iteration gives
\[
\mu^*(A)\geq
\sum_{n=1}^{N}\mu^*(A\cap E_n)
 +\mu^*\!\left(A\setminus\bigcup_{n=1}^{N}E_n\right).
\]
Let $N\to\infty$ and use subadditivity for the reverse inequality. Thus
$\bigcup_nE_n$ is measurable, and arbitrary countable unions follow by
disjointification. The same calculation proves countable additivity of the
restriction.

If $N\in\mathcal M_{\mu^*}$ and $\mu(N)=0$, then every $A\subseteq N$ has
$\mu^*(A)=0$. For every $B\subseteq\Omega$, subadditivity and monotonicity
then give
\[
\mu^*(B)\leq\mu^*(B\cap A)+\mu^*(B\setminus A)
 =\mu^*(B\setminus A)\leq\mu^*(B),
\]
so $A$ is measurable. Hence the measure is complete.

Finally, for $E\in\mathcal R$, cover an arbitrary set $A$ by ring sets
$A_n$. Finite additivity gives
$\mu_0(A_n)=\mu_0(A_n\cap E)+\mu_0(A_n\setminus E)$. Taking infima proves the
Carath\'eodory inequality for $E$. Agreement with $\mu_0$ was proved above.
\end{proof}

\begin{thm}[Carath\'eodory extension and uniqueness]
Every premeasure $\mu_0$ on a ring $\mathcal R$ extends to a measure on
$\sigma(\mathcal R)$. If there are sets $E_k\in\mathcal R$ with
$E_k\uparrow\Omega$ and $\mu_0(E_k)<\infty$, then this extension is unique on
$\sigma(\mathcal R)$.
\end{thm}

\begin{proof}
Existence follows by restricting the induced outer measure to the
Carath\'eodory measurable sets. For uniqueness, let $\mu$ and $\nu$ be two
extensions and fix $k$. The class
\[
\mathcal D_k=
\{A\in\sigma(\mathcal R):\mu(A\cap E_k)=\nu(A\cap E_k)\}
\]
is a Dynkin system containing the $\pi$-system $\mathcal R$, because both
restricted measures have finite total mass on $E_k$. Hence
$\mathcal D_k=\sigma(\mathcal R)$. Letting $k\to\infty$ and using continuity
from below gives $\mu(A)=\nu(A)$ for every $A\in\sigma(\mathcal R)$.
\end{proof}

\subsection{Radon measures}

\begin{prop}[Continuity of a measure]
Let $\mu$ be a measure.
\begin{enumerate}[label=(\roman*)]
  \item If $A_n\uparrow A$, then $\mu(A_n)\uparrow\mu(A)$.
  \item If $A_n\downarrow A$ and $\mu(A_1)<\infty$, then
  $\mu(A_n)\downarrow\mu(A)$.
\end{enumerate}
Conversely, a finitely additive set function on an algebra is a premeasure if
it is continuous from below. If its total mass is finite, this is equivalent
to continuity from above, and also to continuity at the empty set.
\end{prop}

\begin{proof}
For the first assertion, write $A$ as the disjoint union of $A_1$ and the
increments $A_n\setminus A_{n-1}$. The second follows by applying the first to
$A_1\setminus A_n$. The converses follow by applying the corresponding
continuity property to partial unions of a disjoint sequence; under finite
total mass, complements convert increasing sequences into decreasing ones.
\end{proof}

\begin{de}[Radon measure]
A Borel measure $\mu$ on a locally compact Hausdorff space $X$ is a
\emph{Radon measure} if it is finite on compact sets and inner regular:
\[
\mu(B)=\sup\{\mu(K):K\subseteq B,\ K\text{ compact}\}
\qquad(B\in\mathcal B(X)).
\]
On metric spaces this is equivalent, for locally finite Borel measures, to
being both inner regular by compact sets and outer regular by open sets.
\end{de}

\begin{thm}[Regularity on Euclidean space]
Every locally finite Borel measure on $\mathbb R^d$ is Radon. In particular,
every finite Borel measure on $\mathbb R^d$ is inner regular and outer
regular.
\end{thm}

\begin{proof}
First assume that $\mu$ is finite. Let $\mathcal R$ be the class of Borel
sets $B$ such that for every $\varepsilon>0$ there are a compact set
$K\subseteq B$ and an open set $G\supseteq B$ with
$\mu(G\setminus K)<\varepsilon$. Closed sets belong to $\mathcal R$: use the
compact exhaustion $F\cap\overline B(0,n)$ for inner approximation and the
open neighbourhoods $\{x:d(x,F)<1/n\}$ for outer approximation. The class
$\mathcal R$ is a $\sigma$-algebra: complements interchange inner and outer
approximations, while countable unions are handled by approximating a finite
initial union and making the remaining total measure small. Since closed
sets generate the Borel $\sigma$-algebra, every Borel set belongs to
$\mathcal R$.

For a locally finite measure, apply the finite result to the restrictions of
$\mu$ to increasing compact balls and then let the radius tend to infinity.
\end{proof}

\begin{prop}[Distribution function of a locally finite measure]
Let $\mu$ be a locally finite Borel measure on $\mathbb R$. Define
\[
F_\mu(x)=
\begin{cases}
 \mu((0,x]),&x\geq0,\\
 -\mu((x,0]),&x<0.
\end{cases}
\]
Then $F_\mu$ is nondecreasing and right-continuous, and
\[
\mu((a,b])=F_\mu(b)-F_\mu(a)
\qquad(a<b).
\]
Conversely, every nondecreasing right-continuous function determines a unique
locally finite Lebesgue--Stieltjes measure through this identity.
\end{prop}

\begin{proof}
The increment identity follows by splitting intervals at zero and using
finite additivity. Monotonicity is immediate. If $x_n\downarrow x$, then
$(x,x_n]\downarrow\varnothing$ and local finiteness gives
$F_\mu(x_n)-F_\mu(x)\to0$. The converse is the Lebesgue--Stieltjes
construction followed by uniqueness of the $\sigma$-finite extension.
\end{proof}

\begin{thm}[Uniqueness of Lebesgue measure]
If $\mu$ is a translation-invariant, locally finite Borel measure on
$\mathbb R$, then
\[
\mu=c\lambda,
\qquad c=\mu([0,1)),
\]
where $\lambda$ is Lebesgue measure. In particular, the normalization
$\mu([0,1))=1$ determines Lebesgue measure uniquely.
\end{thm}

\begin{proof}
Translation invariance and finite additivity give
$\mu([0,n))=nc$ and $\mu([0,1/n))=c/n$. Hence
$\mu([0,q))=cq$ for every nonnegative rational $q$. Continuity from below and
above extends this identity to every real $q\geq0$. Translation invariance
then gives $\mu([a,b))=c(b-a)$ for all $a<b$. The half-open intervals form a
$\pi$-system generating $\mathcal B(\mathbb R)$, and uniqueness of measures
on that generating class proves $\mu=c\lambda$.
\end{proof}

\subsection{Random variables and measurability}

\begin{de}[Pullback and pushforward]
Let $f:\Omega\to S$ be a map.
\begin{enumerate}[label=(\roman*)]
  \item For a family $\mathcal C\subseteq\mathcal P(S)$, its pullback is
  $f^{-1}(\mathcal C)=\{f^{-1}(C):C\in\mathcal C\}$.
  \item If $(\Omega,\mathcal F,P)$ is a probability space and
  $f:(\Omega,\mathcal F)\to(S,\mathcal S)$ is measurable, the pushforward
  law is
  \[
  f_\#P(B)=P(f\in B)=P(f^{-1}(B)),\qquad B\in\mathcal S.
  \]
\end{enumerate}
\end{de}

\begin{prop}[Generated pullbacks]
For every family $\mathcal C\subseteq\mathcal P(S)$,
\[
f^{-1}(\sigma(\mathcal C))=\sigma(f^{-1}(\mathcal C)).
\]
Consequently, to prove that $f:(\Omega,\mathcal F)\to(S,\sigma(\mathcal C))$
is measurable, it is enough to check $f^{-1}(C)\in\mathcal F$ for
$C\in\mathcal C$.
\end{prop}

\begin{proof}
The class of sets $B\subseteq S$ whose inverse image belongs to
$\sigma(f^{-1}(\mathcal C))$ is a $\sigma$-algebra containing $\mathcal C$.
This gives one inclusion; the other follows because inverse images preserve
complements and countable unions.
\end{proof}

\begin{thm}[Doob--Dynkin factorization]
Let $X:(\Omega,\mathcal F)\to(S,\mathcal S)$ be measurable and let
$Y:\Omega\to\mathbb R$ be $\sigma(X)$-measurable. Then there is an
$\mathcal S/\mathcal B(\mathbb R)$-measurable function $g:S\to\mathbb R$ such
that
\[
Y=g\circ X.
\]
The same conclusion holds for maps into any standard Borel space.
\end{thm}

\begin{proof}
For each rational $q$, choose $A_q\in\mathcal S$ such that
$\{Y<q\}=X^{-1}(A_q)$. Replacing $A_q$ by
$\bigcup_{r<q,\ r\in\mathbb Q}A_r$, we may assume that the family is
increasing and that its inverse images are unchanged. Define the extended-real measurable function
\[
\widetilde g(s)=\inf\{q\in\mathbb Q:s\in A_q\},
\qquad \inf\varnothing=+\infty.
\]
For rational $q$,
$\{\widetilde g<q\}=\bigcup_{r<q}A_r$, so $\widetilde g$ is measurable. Put
$g=\widetilde g$ on $\{\widetilde g\in\mathbb R\}$ and $g=0$ elsewhere.
Then $g:S\to\mathbb R$ is measurable, and the rational sublevel sets show
that $g(X)=Y$. The standard Borel version follows by
transporting the argument through a Borel isomorphism with a Borel subset of
$[0,1]$.
\end{proof}

\begin{thm}[Integration with respect to a law]
Let $X:(\Omega,\mathcal F,P)\to(S,\mathcal S)$ be measurable and let
$\mu_X=X_\#P$. For every nonnegative measurable $h$, and for every integrable
measurable $h$,
\[
\mathbb E[h(X)]=\int_S h(x)\,\mu_X(\dd x).
\]
Moreover, $h(X)$ is integrable if and only if
$\int_S|h|\,\dd\mu_X<\infty$.
\end{thm}

\begin{proof}
The identity is immediate for indicators, extends to nonnegative simple
functions by linearity, to nonnegative functions by monotone convergence, and
to integrable functions by positive and negative parts.
\end{proof}

\begin{lem}[First Borel--Cantelli lemma]
For measurable sets $(A_n)$ in a measure space,
\[
\sum_{n=1}^{\infty}\mu(A_n)<\infty
\quad\Longrightarrow\quad
\mu(\limsup_{n\to\infty}A_n)=0,
\]
where
$\limsup_nA_n=\bigcap_{m\ge1}\bigcup_{n\ge m}A_n$ is the event that infinitely
many $A_n$ occur.
\end{lem}

\begin{proof}
For every $m$,
\[
\mu\!\left(\bigcup_{n\ge m}A_n\right)
 \leq\sum_{n\ge m}\mu(A_n).
\]
The right-hand side tends to zero, and continuity from above gives the
result.
\end{proof}

\begin{thm}[Differentiation under the integral sign]
Let $I\subseteq\mathbb R$ be open, let $(S,\mathcal S,\mu)$ be a measure
space, and let $f:I\times S\to\mathbb R$ satisfy:
\begin{enumerate}[label=(\roman*)]
  \item for every $x\in I$, $f(x,\cdot)$ is measurable and integrable;
  \item for almost every $s$, the map $x\mapsto f(x,s)$ is differentiable;
  \item there is $h\in L^1(\mu)$ such that
  $|\partial_x f(x,s)|\leq h(s)$ for all $x\in I$ and almost every $s$.
\end{enumerate}
Then $F(x)=\int f(x,s)\,\mu(\dd s)$ is differentiable and
\[
F'(x)=\int \partial_xf(x,s)\,\mu(\dd s).
\]
\end{thm}

\begin{proof}
Fix $x\in I$ and let $t_n\to0$ with $x+t_n\in I$. For almost every $s$, the
difference quotient converges to $\partial_xf(x,s)$. The mean-value theorem
bounds its absolute value by $h(s)$. Dominated convergence therefore gives
the derivative along every such sequence, and hence the asserted derivative.
\end{proof}

\begin{de}[Multiplicative class]
A family $\mathcal M$ of bounded real-valued functions on $S$ is
\emph{multiplicative} if $fg\in\mathcal M$ whenever $f,g\in\mathcal M$.
Let $\mathcal H(\mathcal M)$ be the smallest vector space of bounded
functions that contains $1$ and $\mathcal M$ and is closed under bounded
pointwise convergence.
\end{de}

\begin{thm}[Functional monotone-class theorem]
If $\mathcal M$ is multiplicative, then $\mathcal H(\mathcal M)$ contains
every bounded $\sigma(\mathcal M)$-measurable function, where
$\sigma(\mathcal M)$ is the smallest $\sigma$-algebra making every member of
$\mathcal M$ measurable.
\end{thm}

\begin{proof}
For fixed $f\in\mathcal H(\mathcal M)$, let
$\mathcal H_f=\{g:fg\in\mathcal H(\mathcal M)\}$. This class is a vector
space closed under bounded convergence. First take $f\in\mathcal M$; the
multiplicative property shows that $\mathcal M\subseteq\mathcal H_f$, hence
$\mathcal H(\mathcal M)\subseteq\mathcal H_f$. Reversing the roles of $f$
and $g$ shows that $\mathcal H(\mathcal M)$ is itself closed under products.

Now let
$\mathcal D=\{A\subseteq S:\mathbf1_A\in\mathcal H(\mathcal M)\}$. The class
$\mathcal D$ is a Dynkin system. For $f\in\mathcal M$, polynomial functions of
$f$ belong to $\mathcal H(\mathcal M)$ by product closure. Uniform polynomial
approximation on the compact range of $f$ therefore puts every bounded
continuous function of $f$ in $\mathcal H(\mathcal M)$. Approximating the
indicator of $(-\infty,a]$ by bounded continuous cutoff functions shows that
\[
\{f\leq a\}\in\mathcal D
\qquad(f\in\mathcal M,\ a\in\mathbb Q).
\]
Because $\mathcal H(\mathcal M)$ is closed under products, $\mathcal D$ is
also closed under finite intersections. It therefore contains the
$\pi$-system of finite intersections of these sublevel sets, which generates
$\sigma(\mathcal M)$. The $\pi$--$\lambda$ theorem gives
$\sigma(\mathcal M)\subseteq\mathcal D$. Finally, bounded measurable
functions are bounded pointwise limits of simple functions, completing the
proof.
\end{proof}

\begin{prop}
The Borel $\sigma$-algebra on $\mathbb R^d$ is generated by
$C_c(\mathbb R^d)$:
\[
\sigma(C_c(\mathbb R^d))=\mathcal B(\mathbb R^d).
\]
\end{prop}

\begin{proof}
Every continuous function is Borel measurable, so one inclusion is clear.
For the converse, the claim is immediate for $U=\mathbb R^d$. For any other
open set $U$, choose continuous cutoff functions $\eta_m$ that are positive on
$B(0,m)$ and vanish outside $B(0,m+1)$, and put
\[
f_m(x)=\eta_m(x)\min\{1,d(x,U^c)\}.
\]
Then $f_m\in C_c(\mathbb R^d)$ and
$U=\bigcup_m\{f_m>0\}$. Thus every open set belongs to the generated
$\sigma$-algebra.
\end{proof}

\begin{thm}[Determination by compactly supported test functions]
Let $\mu$ and $\nu$ be finite Borel measures on $\mathbb R^d$. If
\[
\int f\,\dd\mu=\int f\,\dd\nu
\qquad\text{for every }f\in C_c(\mathbb R^d),
\]
and $\mu(\mathbb R^d)=\nu(\mathbb R^d)$, then $\mu=\nu$. In particular, the
extra mass condition is automatic for probability measures.
\end{thm}

\begin{proof}
Let $\mathcal M=C_c(\mathbb R^d)\cup\{1\}$; this is multiplicative. The class
of bounded Borel functions $f$ for which the two integrals agree is a vector
space containing $\mathcal M$ and closed under bounded pointwise convergence.
The functional monotone-class theorem and the preceding proposition imply
that it contains every bounded Borel function. Apply this to indicators.
\end{proof}

\subsection{Product measures}

\begin{de}[Product $\sigma$-algebra]
For measurable spaces $(S_i,\mathcal S_i)$, $i\in I$, the product
$\sigma$-algebra on $\prod_{i\in I}S_i$ is
\[
\bigotimes_{i\in I}\mathcal S_i
 :=\sigma\{\pi_i^{-1}(A):i\in I,\ A\in\mathcal S_i\},
\]
where $\pi_i$ is the $i$th coordinate projection.
\end{de}

\begin{prop}[Coordinate criterion]
A map $X:(\Omega,\mathcal F)\to
(\prod_{i\in I}S_i,\bigotimes_{i\in I}\mathcal S_i)$ is measurable if and
only if every coordinate $\pi_i\circ X$ is measurable.
\end{prop}

\begin{proof}
Necessity follows because each projection is measurable. Conversely, the
inverse image under $X$ of every generating cylinder
$\pi_i^{-1}(A)$ is $(\pi_i\circ X)^{-1}(A)\in\mathcal F$.
\end{proof}

\begin{thm}[Product measure and Tonelli--Fubini]
Let $(S,\mathcal S,\mu)$ and $(T,\mathcal T,\nu)$ be $\sigma$-finite measure
spaces. There is a unique measure $\mu\otimes\nu$ on
$\mathcal S\otimes\mathcal T$ satisfying
\[
(\mu\otimes\nu)(A\times B)=\mu(A)\nu(B).
\]
If $f:S\times T\to[0,\infty]$ is measurable, then its sections are
measurable and
\[
\int f\,\dd(\mu\otimes\nu)
 =\int_S\!\left(\int_T f(s,t)\,\nu(\dd t)\right)\mu(\dd s)
 =\int_T\!\left(\int_S f(s,t)\,\mu(\dd s)\right)\nu(\dd t).
\]
If $f$ is integrable, the same identities hold for $f$, both iterated
integrals are absolutely integrable, and the order of integration may be
interchanged.
\end{thm}

\begin{proof}
Construct the product measure from the rectangle premeasure. The identities
hold first for rectangle indicators, then for nonnegative simple functions,
and then for nonnegative measurable functions by monotone convergence. For
integrable $f$, apply the nonnegative result to $f^+$ and $f^-$, or first to
$|f|$, to justify all subtractions and both iterated integrals.
\end{proof}

\subsection{Kolmogorov's extension theorem}

\begin{thm}[Kolmogorov extension on the compact cube]
Let $Q=[0,1]^{\mathbb N}$. For each $n$, let $\mu_n$ be a probability measure
on $[0,1]^n$. Assume projective consistency:
\[
\mu_{n+1}(B\times[0,1])=\mu_n(B)
\qquad(B\in\mathcal B([0,1]^n)).
\]
Then there is a unique probability measure $\mu$ on
$(Q,\bigotimes_{n\ge1}\mathcal B([0,1]))$ such that
\[
\mu(\pi_{1:n}^{-1}(B))=\mu_n(B)
\qquad(n\ge1).
\]
\end{thm}

\begin{proof}
Define $\mu_0$ on the cylinder algebra by
$\mu_0(\pi_{1:n}^{-1}(B))=\mu_n(B)$. Consistency makes this well defined and
finitely additive. It remains to prove continuity at the empty set. Suppose
$C_k$ is a decreasing sequence of cylinders with
$\inf_k\mu_0(C_k)>0$. By regularity of the finite-dimensional measures,
choose compact cylinders $K_k\subseteq C_k$ so that
$\mu_0(C_k\setminus K_k)<2^{-k-1}\inf_j\mu_0(C_j)$. Since
$C_m\subseteq C_k$ for $k\leq m$,
\begin{align*}
\mu_0\!\left(C_m\setminus\bigcap_{k=1}^{m}K_k\right)
&\leq\sum_{k=1}^{m}\mu_0(C_m\setminus K_k)\\
&\leq\sum_{k=1}^{m}\mu_0(C_k\setminus K_k)
 <\frac12\inf_j\mu_0(C_j).
\end{align*}
Thus every finite intersection $\bigcap_{k=1}^{m}K_k$ has positive measure
and is nonempty. These intersections are compact and decreasing in the compact
space $Q$, so their full intersection is nonempty. Hence $\bigcap_kC_k\neq\varnothing$. The
contrapositive is continuity at the empty set, so $\mu_0$ is a premeasure.
Carath\'eodory's theorem gives existence, and uniqueness follows because
cylinders form a generating $\pi$-system.
\end{proof}

\begin{de}[Standard Borel space]
A measurable space $(S,\mathcal S)$ is \emph{standard Borel} if it is
measurably isomorphic to a Borel subset of a Polish space. Equivalently, it is
measurably isomorphic to a Borel subset of $[0,1]$.
\end{de}

\begin{thm}[Kolmogorov extension for countable standard Borel products]
Let $(S_n,\mathcal S_n)$ be standard Borel spaces. Suppose $\mu_n$ is a
probability measure on
$(S_1\times\cdots\times S_n,\mathcal S_1\otimes\cdots\otimes\mathcal S_n)$
and the family is projectively consistent. Then there is a unique probability
measure $\mu$ on
\[
\left(\prod_{n=1}^{\infty}S_n,
      \bigotimes_{n=1}^{\infty}\mathcal S_n\right)
\]
with finite-dimensional marginals $\mu_n$.
\end{thm}

\begin{proof}
Choose Borel isomorphisms from $S_n$ onto Borel subsets $B_n\subseteq[0,1]$.
Transport each finite-dimensional law to $[0,1]^n$, extending it by zero
outside $B_1\times\cdots\times B_n$. The compact-cube theorem gives a measure
on $[0,1]^{\mathbb N}$. Every coordinate belongs to its corresponding $B_n$
with probability one, so the countable product $\prod_nB_n$ has full measure.
Restrict to this set and transport the measure back. Uniqueness again follows
from the cylinder $\pi$-system.
\end{proof}

\begin{thm}[Infinite product measures]
Let $(\nu_n)_{n\ge1}$ be Borel probability measures on $\mathbb R$. There is
a unique probability measure
$\nu=\bigotimes_{n=1}^{\infty}\nu_n$ on
$(\mathbb R^{\mathbb N},\bigotimes_{n\ge1}\mathcal B(\mathbb R))$ such that
\[
\nu\!\left(A_1\times\cdots\times A_n\times
           \prod_{k>n}\mathbb R\right)
 =\prod_{k=1}^{n}\nu_k(A_k).
\]
For every nonnegative measurable, or integrable, function
$f:\mathbb R^n\to\mathbb R$,
\[
\int_{\mathbb R^{\mathbb N}}f(x_1,\ldots,x_n)\,\nu(\dd x)
 =\int_{\mathbb R^n}f(x_1,\ldots,x_n)
   \prod_{k=1}^{n}\nu_k(\dd x_k).
\]
\end{thm}

\begin{proof}
Apply the preceding theorem to the consistent finite-dimensional measures
$\nu_1\otimes\cdots\otimes\nu_n$. The integral identity is the
finite-dimensional pushforward formula.
\end{proof}

\begin{cor}[Construction of an independent sequence]
Let $(\mu_n)$ be Borel probability measures on $\mathbb R$. There is a
probability space carrying independent random variables $(X_n)$ such that
$\mathcal L(X_n)=\mu_n$ for every $n$. They are identically distributed
precisely when all the measures $\mu_n$ are equal.
\end{cor}

\begin{proof}
Take $\Omega=\mathbb R^{\mathbb N}$ with the product measure
$\bigotimes_n\mu_n$, and let $X_n=\pi_n$. The cylinder formula gives both the
marginal laws and independence.
\end{proof}
